\documentclass[12pt]{article}
\usepackage[utf8]{inputenc}
\usepackage[T1]{fontenc}
\usepackage[margin=1in]{geometry}
\usepackage{setspace}
\usepackage{xcolor}
\usepackage[colorlinks=true, urlcolor=blue!60!black, linkcolor=blue!60!black]{hyperref}
\setlength{\parindent}{2em}
\setlength{\parskip}{1.25em}
\renewcommand{\baselinestretch}{1.0}
\begin{document}
\begin{center}
{\LARGE \textit{Ward v. Rock Against Racism}, 491 U.S. 781 (1989)\par}
\vspace{0.8em}
{\large Case Brief --- Four-Page Edition\par}
\end{center}
\vspace{1em}

\section{Memory Jogger}

Upholding New York City's rule that Central Park bandshell performers use city sound equipment and a city technician, the Court held that a content-neutral time, place, and manner regulation must serve a significant governmental interest, burden no substantially greater amount of speech than necessary, and leave open ample alternative channels of communication. It need not be the least restrictive means (\href{https://supreme.justia.com/cases/federal/us/491/781/}{Justia}).

\section{Detailed Case Facts}

Rock Against Racism (RAR), an association promoting antiracist views, sponsored annual programs of speeches and rock music from 1979 to 1986 at the Naumburg Acoustic Bandshell in Central Park. The bandshell faces the Sheep Meadow, a designated quiet zone, and apartments on Central Park West. RAR's concerts drew many noise complaints; at its 1984 concert, performers ignored warnings and two citations, and when police cut the power the crowd became abusive. Other sponsors with poor equipment or unskilled technicians produced inadequate sound, frustrating crowds. After rejecting fixed decibel caps (sound varies with temperature, foliage, and crowd size), the city adopted 1986 Use Guidelines requiring all performers to use city-leased, high-quality equipment run by an independent, experienced technician retained by the city.

\section{Procedural History}

RAR brought a facial First Amendment challenge in the Southern District of New York. After a five-day trial, the court upheld the guideline, finding that the city's system met RAR's needs, that sponsors controlled the sound ``mix,'' and that officials consulted sponsors before lowering volume. The Second Circuit reversed, requiring the ``least intrusive'' means and pointing to volume limiters or pulling the plug. The Supreme Court reversed on June 22, 1989.

\section{Judicial Votes}

Reversed, \textbf{6--3}.
\begin{itemize}
\item \textbf{Majority (5):} Justice Kennedy, joined by Chief Justice Rehnquist and Justices White, O'Connor, and Scalia.
\item \textbf{Concurrence:} Justice Blackmun, in the result, without opinion.
\item \textbf{Dissent:} Justice Marshall, joined by Justices Brennan and Stevens.
\end{itemize}

\section{Holding}

Music is protected expression, and the bandshell is a traditional public forum, but the sound-amplification guideline is a valid content-neutral time, place, and manner regulation. Narrow tailoring under intermediate scrutiny does not require the ``least restrictive or least intrusive means''; a regulation is narrowly tailored if it promotes a substantial interest that would be achieved less effectively without it and does not burden substantially more speech than necessary.

\section{Analysis of Opinions}

\textbf{Majority (Kennedy).} Applying the \textit{Clark v. CCNV} framework, Kennedy found the rule content-neutral: the government's purpose controls, and controlling noise to protect the Sheep Meadow and nearby residences had nothing to do with message. The interest in sound quality did not make it content-based, since the city disclaimed aesthetic standards and its technician deferred to sponsors on mix. Unlike the licensing scheme in \textit{City of Lakewood}, the guideline's text and the city's narrowing construction confined discretion, so it was no prior restraint. The city had substantial interests in protecting citizens from unwelcome noise and ensuring adequate sound. The Second Circuit wrongly imported a least-intrusive-means test; that inquiry ``has never been a part of the inquiry into the validity of a time, place, and manner regulation.'' Courts defer to the decisionmaker's reasonable choice and judge fit against the ``overall problem.'' Ample alternatives remained because no type of music was banned and only volume was moderated.

\textbf{Blackmun.} Concurred in the result without opinion.

\textbf{Marshall, with Brennan and Stevens (dissenting).} The guideline was content-neutral and served a significant interest, but the majority ``eviscerat[ed]'' narrow tailoring and substituted ``mandatory deference.'' Under \textit{Schneider} and \textit{Martin v. Struthers}, courts must weigh efficacy against burden when less intrusive tools, such as enforcing noise ordinances or decibel monitoring, exist; now any rule advancing a goal ``in the slightest'' survives. City control of equipment is a prior restraint, like Crown control of printing presses, without standards or prompt review under \textit{Freedman v. Maryland}.

\textbf{Later developments.} \textit{McCullen v. Coakley} (2014) reiterated that \textit{Ward} does not demand the least restrictive means but held that government ``may not sacrifice speech for efficiency.'' \textit{Reed} (2015) made facial content distinctions trigger strict scrutiny regardless of purpose; \textit{City of Austin} (2022) tempered \textit{Reed}.

\section{Examples: Future Applications}

\begin{enumerate}
\item \textbf{Same-Side: Amphitheater by a Hospital.} A city requires events at a waterfront amphitheater beside a hospital to use a city sound-limiting system; performers choose content and mix. The rule is content-neutral, serves a substantial interest, and leaves ample alternatives.
\item \textbf{Same-Side: Courthouse Plaza Hours.} A county bars amplified sound on its courthouse plaza during weekday court hours. That a deputy-enforced decibel limit would work nearly as well does not matter; the rule ``need not be the least restrictive or least intrusive means.''
\item \textbf{Opposite-Side: Standardless Mute Button.} A city requires evening plaza speakers to use a city PA whose technician may mute speech that ``unreasonably disturbs'' commerce or tranquility. With no narrowing construction, it invites an unbridled-discretion attack, and \textit{McCullen} demands a showing that less intrusive tools would fail.
\item \textbf{Opposite-Side: Volume Keyed to Lyrics.} A technician may lower the volume of ``offensive or hateful lyrics.'' The rule is content-based and must survive strict scrutiny.
\item \textbf{Fence-Sitter: Total Ban on the Main Lawn.} A city bans amplification on the traditional rally lawn, allowing it only at a remote stage. The rule is neutral, but unlike \textit{Ward} it may leave inadequate alternatives for large rallies.
\end{enumerate}

\section{Critique}

\textbf{\textit{Ward} versus \textit{McCullen}.} \textit{Ward} deferred to New York's rejection of decibel caps, while \textit{McCullen} required Massachusetts to show that less intrusive tools failed. \textit{McCullen} may quietly revive Justice Marshall's balancing whenever a rule heavily burdens core one-on-one or sidewalk political speech.

\textbf{Unwritten narrowing constructions.} \textit{Ward} upheld broad text (``best sound'') partly because officials testified they deferred to performers. Critics say reliance on such practice weakens the facial safeguard against unbridled discretion when administrations change.

\section{Key Quotations}

\begin{quote}
``Lest any confusion on the point remain, we reaffirm today that a regulation of the time, place, or manner of protected speech must be narrowly tailored to serve the government's legitimate, content-neutral interests, but that it need not be the least restrictive or least intrusive means of doing so.''
\end{quote}
\noindent --- Justice Kennedy, for the Court, 491 U.S. at 798--799.

\begin{quote}
``The principal inquiry in determining content-neutrality, in speech cases generally and in time, place, or manner cases in particular, is whether the government has adopted a regulation of speech because of disagreement with the message it conveys. \dots{} The government's purpose is the controlling consideration.''
\end{quote}
\noindent --- Justice Kennedy, for the Court, 491 U.S. at 791.

\begin{quote}
``Music is one of the oldest forms of human expression. \dots\ Music, as a form of expression and communication, is protected under the First Amendment.''
\end{quote}
\noindent --- Justice Kennedy, for the Court, 491 U.S. at 790.

\end{document}
